\section{Products in the continuous coefficient curvature module}
\label{sec:r6-continuous-curvature-module-products}

This section records a consequence of the represented-order theorem.
The input currents have a prescribed local finite potential
presentation. No product of completely arbitrary positive closed
currents is defined.

\begin{definition}
On an allowed ordinary open $U$, let $\mathscr M^1(U)$ be the
$C(U,\mathbf R)$-module of original coefficient-Radon $(1,1)$-currents
locally admitting a finite expression
\[
                         A=\sum_{\alpha}f_\alpha D[g_\alpha],
 \qquad f_\alpha\in C(U,\mathbf R),\quad g_\alpha\in\mathcal C_b(U).
\]
The expression is local; actual equality of currents is the equivalence
relation. Signed bounded classical potentials can be used after
expanding them into two classical parts. Its positive cone consists
of those represented currents which are positive in the original
middle CLD sense. Variable continuous coefficients need not preserve
separate closedness.
\end{definition}

\begin{theorem}[Positive products of continuously weighted curvatures]
\label{thm:r6-positive-continuous-module-products}
Let $S$ be any fixed allowed $(p,p)$-background and $r\le n-p$.
There is a well-defined symmetric $C(U,\mathbf R)$-multilinear map
\[
 \Phi_{S,r}:(\mathscr M^1)^r\longrightarrow
       \{\text{original coefficient-Radon $(p+r,p+r)$-currents}\}
\]
characterized on finite local presentations by
\begin{equation}
 \Phi_{S,r}(A_1,\ldots,A_r)
 =\sum_{\alpha_1,\ldots,\alpha_r}
    \left(\prod_{i=1}^r f_{i,\alpha_i}\right)
    T_{S,r}(g_{1,\alpha_1},\ldots,g_{r,\alpha_r}).
 \label{eq:r6-continuous-module-product}
\end{equation}
It commutes with ordinary restriction and gluing. If every $A_i$ is
positive, its image is positive. In particular the map is increasing
in each slot when all the other slots are positive. When every slot
is $D[g_i]$, it is the frozen bounded BT product.

No separate closedness of the output is asserted for general
continuous coefficients. The map does not factor through every
higher-degree ordinary-current relation merely from its construction.
\end{theorem}

\noindent\textbf{Status: \texttt{DIRECT\_COROLLARY}} of the new
bounded-difference positivity criterion and continuous-coefficient
relation theorem, with the positivity argument written out.

\begin{proof}
Each expression in \eqref{eq:r6-continuous-module-product} is a finite
sum of continuous multiples of positive coefficient-Radon currents.
Suppose the presentation of one slot is changed without changing its
original current. Fix one multi-index in the remaining slots and put
\[
 R=T_{S,r-1}(g_{1,\alpha_1},\ldots,
              \widehat{g_{i,\alpha_i}},\ldots,g_{r,\alpha_r}).
\]
This is an allowed positive separately closed background. The
zero-relation assertion of
Theorem~\ref{thm:r6-bounded-curvature-continuous-module}, applied to
$R$, makes the change in that slot zero. Multiplication by the
remaining continuous coefficient product and finite summation
preserve zero. Changing the slots one at a time proves presentation
independence. Finite linearity, scalar multiplication, symmetry of
the frozen BT tails, ordinary locality and gluing give the stated
module and sheaf properties.

For positivity fix a positive original compact complementary test
$\eta$. Use finitely many local presentations near its support and
nonnegative smooth localization; it suffices to treat one presentation.
All its coefficients are bounded by a fixed number $M\ge1$ on a
slightly larger compact neighborhood. For $0<\epsilon<1$, subdivide
the test support into finitely many ordinary neighborhoods on which
\[
                    |f_{i,\alpha}-c_{i,\alpha}|<\epsilon
\]
for suitable constants $c_{i,\alpha}$. On one such neighborhood put
\[
 H_i=\sum_\alpha(c_{i,\alpha}+\epsilon)g_{i,\alpha}.
\]
It is an actual bounded classical difference, and
\[
 D[H_i]=A_i+\sum_\alpha(c_{i,\alpha}+\epsilon-f_{i,\alpha})
                                      D[g_{i,\alpha}]\ge0.
\]
Theorem~\ref{thm:r6-positive-bounded-dsh-classical} therefore gives
$H_i\in\mathcal C_b$ on that neighborhood. Hence
$T_{S,r}(H_1,\ldots,H_r)\ge0$ there.

For every multi-index the scalar difference between its coefficient
in \eqref{eq:r6-continuous-module-product} and its coefficient in
$T_{S,r}(H_1,\ldots,H_r)$ has absolute value at most
$2r\epsilon(M+2)^{r-1}$. This follows by telescoping the product of
$r$ scalars, since each paired scalar differs by at most $2\epsilon$.
Let $R_0$ be the finite sum of all the positive BT tails in
\eqref{eq:r6-continuous-module-product}, without coefficients.
On this neighborhood the positive test cone therefore gives
\[
          \Phi_{S,r}(A_1,\ldots,A_r)
                   \ge -2r\epsilon(M+2)^{r-1}R_0.
\]
The same constant and the same $R_0$ work for every member of the
finite subdivision. Apply a nonnegative original smooth partition
to $\eta$, sum the inequalities, and let $\epsilon\downarrow0$.
The fixed control mass $\langle R_0,\eta\rangle$ is finite. The
pairing of $\Phi_{S,r}$ with $\eta$ is nonnegative. This proves
positivity on all original tests in the specified cone.

For monotonicity subtract the two products and telescope their
slots. Each term has one positive difference slot and all other
slots positive, so the proved positivity applies. The BT
specialization is its defining one-term presentation.
\end{proof}

\begin{example}[Borel coefficients do not preserve current relations]
\label{ex:r6-borel-curvature-module-failure}
The word continuous in the preceding module cannot be replaced by
bounded Borel, even when all potential functions are continuous.
Take the ordinary bidannulus
\[
 U=\{r<|T_1|<R,\ r<|T_2|<R\}\subset(\mathbf G_{m,K}^2)^{\an},
             \qquad 0<r<1<R,
\]
and let $Y=\{T_2=1\}\cap U$, $S=[Y]$, and
$g=\max\{0,\log|T_1|\}$. This is a continuous locally bounded
classical psh function. The frozen first-order analytic-divisor
nonconcentration theorem, applied to $D[g]$, gives
\[
                         \mathbf1_YD[g]=0
\]
as an ambient coefficient-Radon current. On the other hand, the
restriction and projection formulas for the subvariety background,
and the one-variable slope computation, give
\[
 \mathbf1_YD(g[Y])=D(g[Y])=i_*\delta_{\zeta_{0,1}}\ne0,
                  \qquad i:Y\hookrightarrow U.
\]
Here $\zeta_{0,1}$ is the radius-one Gauss point in the $T_1$ annulus;
the slope jump of $\max(0,\log|T_1|)$ is exactly one. An original
nonnegative compact smooth scalar test equal to one near this point
pairs with the displayed measure to give one. Thus the failure is
visible on original tests, not only on an enlarged test sheaf.
\end{example}

\noindent\textbf{Status: \texttt{FALSE}} for the proposed Borel
extension; the example is a \texttt{DIRECT\_COROLLARY} of the frozen
polar-carrier obstruction and first-order nonconcentration. It is not
counted as a new counterexample mechanism. R5,
Corollary~\ref{cor:r5-universal-borel-kernel-strict} and
Remark~\ref{rem:r5-bounded-usc-kernel-obstruction}, already show why
arbitrary assigned bounded upper semicontinuous values must also be
excluded from the bounded-classical-difference positivity criterion.

\paragraph{What this construction does not settle.}
Independence of the degree-one presentations means that relations in
$\mathscr M^1$ may be substituted inside
\eqref{eq:r6-continuous-module-product}. It does not say that every
ordinary zero current formed from several curvature factors is such
a relation. The higher-degree analytic delta kernel remains a separate
problem. Nor is the output of a continuous coefficient product an
admissible new background unless its required positivity and separate
closedness are independently available.
